\chapter{PHÂN TÍCH VÀ THIẾT KẾ MÔ HÌNH}
\label{chap:chapter3}

\section{Mô hình hóa bài toán tồn kho đa kho}
\label{sec:moHinhHoa}

\subsection{Mô tả mạng lưới}
\label{subsec:mangLuoi}

Hệ thống gồm $N = 10$ kho hoạt động song song và độc lập về nguồn cung, tương
ứng với $10$ cửa hàng của bộ dữ liệu M5 Walmart. Mỗi kho lưu trữ $M = 30$ mặt
hàng (SKU) đại diện, được chọn theo phương pháp phân tầng trên các nhóm ngành
hàng và lọc bỏ những mặt hàng có nhu cầu quá thấp (dưới $0{,}2$ đơn vị/ngày ở
đa số cửa hàng) để tránh đưa vào hệ thống những cặp gần như không phát sinh
giao dịch, vốn chỉ làm nhiễu tín hiệu học của mạng chia sẻ tham số. Tổng cộng
hệ thống có $N \times M = 300$ cặp kho--mặt hàng, mỗi cặp là một tác tử ra
quyết định độc lập trên quan sát cục bộ của chính nó.

Mỗi kho $k$ có dung lượng hữu hạn $W_k$, được hiệu chỉnh \emph{riêng theo nhu
cầu của chính kho đó} thay vì dùng một con số chung cho toàn hệ thống:
\begin{equation}
    W_k = c_{\text{cover}} \sum_{i \in k} \bar{D}_i
    \label{eq:capacity}
\end{equation}
trong đó $\bar{D}_i$ là nhu cầu trung bình hàng ngày của cặp $i$ (ước lượng chỉ
từ giai đoạn dữ liệu dùng để huấn luyện, xem Mục~\ref{subsec:phanChiaMien}) và
$c_{\text{cover}} = 5$ là số ngày dự trữ mục tiêu. Việc hiệu chỉnh theo từng
kho là cần thiết vì quy mô kinh doanh giữa $10$ cửa hàng M5 chênh lệch nhau
nhiều lần: nếu dùng một dung lượng chung cho tất cả, kho có doanh số lớn sẽ chỉ
đủ chứa cho chưa tới hai ngày nhu cầu -- thấp hơn cả thời gian giao hàng tối đa
-- khiến mục tiêu mức phục vụ trở thành bài toán vô nghiệm với kho đó, trong
khi kho nhỏ lại dư thừa dung lượng đến mức ràng buộc sức chứa không còn tác
dụng ghép nối các mặt hàng với nhau.

Sự ràng buộc giữa các mặt hàng trong cùng kho phát sinh chính từ dung lượng
dùng chung $W_k$ này: khi một mặt hàng chiếm nhiều chỗ, các mặt hàng khác
trong cùng kho có nguy cơ bị từ chối nhập hàng vượt mức. Đây là yếu tố khiến
bài toán mang bản chất hợp tác và không thể tách rời thành $300$ bài toán đơn
lẻ độc lập.

Nhu cầu hàng ngày $D_t^{(i)}$ cho mỗi cặp kho--mặt hàng $i$ được lấy từ hai
nguồn:
\begin{itemize}
    \item Dữ liệu M5 Walmart: chuỗi doanh số thực tế của $10$ cửa hàng, mỗi
    cửa hàng $30$ SKU đại diện đã qua chọn lọc \parencite{makridakis2022background}.
    \item Dữ liệu giả lập: nhu cầu Poisson $D_t \sim
    \text{Poisson}(\lambda_i)$ với $\lambda_i$ khác nhau giữa các SKU, phục vụ
    kiểm tra đơn vị và thử nghiệm đối chứng trong điều kiện kiểm soát.
\end{itemize}

\subsection{Thiết lập bài toán dưới dạng MDP đa tác tử}
\label{subsec:mdpDaTacTu}

Bài toán được mô hình hóa dưới dạng Dec-POMDP \parencite{oliehoek2016} với
$300$ tác tử. Mỗi tác tử $i$ (tương ứng cặp kho--mặt hàng) có:

Không gian quan sát $o_i^t \in [0, 1]^d$: vector quan sát cục bộ tại
kỳ $t$, đã được chuẩn hóa và giới hạn trong đoạn $[0, 1]$ để ổn định quá trình
huấn luyện mạng nơ-ron, bao gồm:
\begin{itemize}
    \item Mức tồn kho hiện tại, quy đổi theo số ngày cầu điển hình của chính
    cặp đó.
    \item Vị thế tồn kho (tồn kho hiện có cộng lượng hàng đang trên đường vận
    chuyển), cùng đơn vị quy đổi.
    \item Lịch sử nhu cầu $7$ ngày gần nhất, biểu diễn theo thang log để giữ
    được thông tin ở cả những ngày cầu đột biến lẫn những ngày cầu bằng
    không -- đặc trưng phổ biến của nhu cầu bán lẻ gián đoạn.
    \item Lượng hàng đang trên đường ứng với từng thời điểm còn lại cho tới
    khi giao hàng.
    \item Tỷ lệ đáp ứng nhu cầu trên cửa sổ trượt $w = 30$ ngày,
    $\text{FR}_t^{(i, w)}$ -- thành phần bắt buộc phải có trong quan sát vì
    hàm phần thưởng (Mục~\ref{sec:hamThuong}) phạt trực tiếp theo đại lượng
    này; nếu thiếu, tác tử không thể suy ra được vì sao mình bị phạt và môi
    trường mất tính Markov tại đúng thành phần đó.
    \item Một đặc trưng quy mô cầu (theo thang log), giúp mạng dùng chung phân
    biệt được một SKU đang quan sát là loại bán chậm hay bán nhanh.
    \item Mức sử dụng sức chứa hiện tại của kho chứa cặp đó, $\sum_{j \in k}
    I_t^{(j)} / W_k$ -- kênh duy nhất để một tác tử ``cảm nhận'' áp lực từ các
    mặt hàng khác cùng kho, và là điều kiện bắt buộc để môi trường giữ tính
    Markov tại thành phần phạt tràn kho (vốn tính trên tổng tồn kho của cả
    kho chứ không phải của riêng từng cặp).
    \item Tỷ trọng tồn kho của chính cặp đó trong tổng tồn kho của kho.
    \item Khi sử dụng dữ liệu giá bán thực tế: một tín hiệu phản ánh mức độ
    giảm giá hiện tại so với giá bán trung bình của chính cặp đó -- một chỉ
    báo sớm cho khả năng nhu cầu tăng đột biến do khuyến mãi.
    \item Ngày trong tuần dạng one-hot ($7$ chiều) và, khi dùng dữ liệu M5
    thực tế, các đặc trưng lịch (sự kiện lễ hội, chương trình trợ cấp thực
    phẩm SNAP, tháng trong năm).
\end{itemize}

Không gian hành động $a_i^t \in \{0, 1, \ldots, K-1\}$: mỗi tác tử chọn
một trong $K = 6$ mức đặt hàng. Vì quy mô nhu cầu giữa các cặp kho--mặt hàng
chênh lệch tới vài trăm lần, các mức đặt hàng được thiết kế \emph{tương đối
theo nhu cầu của chính từng cặp} thay vì theo số lượng tuyệt đối dùng chung:
\begin{equation}
    q_i(a) = \Big\lceil m_a \cdot \bar{D}_i \cdot \bar{L} \Big\rceil,
    \qquad m_a \in \{0;\, 0{,}5;\, 1;\, 2;\, 3;\, 5\}
    \label{eq:orderTable}
\end{equation}
trong đó $\bar{L}$ là thời gian giao hàng trung bình. Bảng lượng đặt hàng
$q_i(\cdot)$ được ép tăng nghiêm ngặt giữa các mức liên tiếp (mức sau luôn lớn
hơn mức trước ít nhất một đơn vị, trừ mức $0$ luôn cố định là không đặt hàng);
nếu không ép điều kiện này, làm tròn một hệ số nhân nhỏ trên một nhu cầu nhỏ có
thể khiến nhiều mức đặt hàng khác nhau quy về cùng một số lượng thực tế, khiến
tác tử nhận được cùng một hệ quả cho các hành động khác nhau và triệt tiêu tín
hiệu gradient cần thiết để chính sách hội tụ. Nhờ thiết kế tương đối này, một
mức đặt hàng ``trung bình'' luôn tương ứng với một số ngày dự trữ có ý nghĩa
như nhau đối với cả SKU bán chậm lẫn SKU bán chạy nhất, dù không gian hành
động mỗi tác tử chỉ có kích thước cố định $K = 6$. Không gian hành động chung
của toàn hệ thống vẫn là tích Descartes $K^{300} = 6^{300}$.

Hàm chuyển trạng thái diễn ra theo trình tự sau tại mỗi kỳ:
\begin{enumerate}
    \item Tác tử quan sát $o_i^t$ và chọn hành động $a_i^t$, quy đổi thành
    lượng đặt hàng $q_i(a_i^t)$ theo công thức~(\ref{eq:orderTable}).
    \item Lô hàng đã đặt trước đó và đã hết thời gian giao hàng được nhập vào
    kho.
    \item Lượng hàng vượt quá dung lượng còn lại của kho bị từ chối ngay tại
    thời điểm nhập kho (không phải bị loại bỏ sau khi đã tính vào tồn kho),
    phân bổ theo đúng tỷ trọng tồn kho hiện có của từng mặt hàng trong kho đó.
    \item Nhu cầu $D_t^{(i)}$ phát sinh; hàng được bán tối đa bằng tồn kho
    hiện có, phần nhu cầu vượt quá bị mất hoàn toàn (mô hình mất đơn hàng).
    \item Tỷ lệ đáp ứng trên cửa sổ trượt và chi phí, phần thưởng của kỳ được
    cập nhật.
    \item Đơn hàng $q_i(a_i^t)$ được ghi nhận vào hàng đợi giao hàng ứng với
    một thời gian giao hàng ngẫu nhiên $L \sim \text{Uniform}(1, 3)$ ngày, hệ
    thống chuyển sang kỳ $t + 1$.
\end{enumerate}

\subsection{Phân chia miền dữ liệu huấn luyện -- kiểm định -- kiểm tra}
\label{subsec:phanChiaMien}

Để tránh rò rỉ thông tin giữa giai đoạn huấn luyện và giai đoạn báo cáo kết
quả, chuỗi thời gian $1.941$ ngày của M5 được cắt thành ba miền liên tiếp và
không chồng lấn: miền huấn luyện (train), miền kiểm định (validation) và miền
kiểm tra (test). Mọi tham số ước lượng từ dữ liệu -- nhu cầu trung bình và độ
lệch chuẩn của từng cặp, dung lượng mỗi kho theo công thức~(\ref{eq:capacity}),
bảng lượng đặt hàng theo công thức~(\ref{eq:orderTable}), và đơn giá bán trung
bình dùng cho chi phí thiếu hàng ở Mục~\ref{subsec:thamSoChiPhi} -- đều chỉ
được tính trên miền huấn luyện. Mỗi tập mô phỏng (episode) dài $365$ ngày luôn
được giới hạn để không vượt ra khỏi ranh giới của miền mà nó thuộc về.

Miền kiểm định được dùng riêng để chọn bộ trọng số tốt nhất trong quá trình
huấn luyện (thông qua đánh giá định kỳ theo chính sách xác định, không lấy
mẫu ngẫu nhiên), còn miền kiểm tra chỉ được truy cập đúng một lần, ở bước đánh
giá cuối cùng dùng để báo cáo và so sánh với các chính sách đối chứng
(Chương~\ref{chap:chapter4}). Việc tách bạch hai miền này là cần thiết: nếu bộ
trọng số tốt nhất được chọn bằng chính hiệu năng đo trên miền dùng để báo cáo
kết quả sau cùng, số liệu công bố sẽ mang một dạng rò rỉ nhẹ, vì bước chọn mô
hình đã ``nhìn thấy'' trước đúng dữ liệu mà nó sẽ được đánh giá trên đó.

\section{Thiết kế hàm phần thưởng}
\label{sec:hamThuong}

\subsection{Phần thưởng phân rã cục bộ}
\label{subsec:localReward}

Mỗi tác tử $i$ nhận phần thưởng riêng dựa trên chi phí của chính cặp kho--mặt
hàng mà nó quản lý, thay vì dùng chung phần thưởng toàn cục:
\begin{equation}
    r_i^t = -\left[
    c_{lk} \cdot I_t^{(i)} +
    c_{th}^{(i)} \cdot U_t^{(i)} +
    c_{dh} \cdot \mathbb{1}(a_i^t > 0) +
    p_{tk} \cdot O_t^{(i)}
    + \phi_{dv} \cdot c_{th} \cdot \bar{D}_i \cdot
      \max\bigl(\tau - \text{FR}_t^{(i, w)},\; 0\bigr)
    \right]
    \label{eq:localReward}
\end{equation}
trong đó:
\begin{itemize}
    \item $I_t^{(i)}$: mức tồn kho sau khi bán hàng trong kỳ $t$.
    \item $U_t^{(i)} = \max(D_t^{(i)} - I_t^{(i)}, 0)$: lượng nhu cầu không
    được đáp ứng.
    \item $O_t^{(i)}$: lượng hàng bị từ chối do vượt dung lượng kho.
    \item $c_{th}^{(i)}$: đơn giá thiếu hàng của cặp $i$ -- có thể là hằng số
    $c_{th}$ dùng chung, hoặc suy ra từ giá bán thực tế của chính mặt hàng đó
    (Mục~\ref{subsec:thamSoChiPhi}).
    \item $\text{FR}_t^{(i, w)}$: tỷ lệ đáp ứng trên cửa sổ trượt $w = 30$
    ngày; $\tau = 0{,}85$: ngưỡng mức phục vụ mục tiêu.
\end{itemize}

Thành phần phạt mức phục vụ được nhân thêm với $\bar{D}_i$ -- nhu cầu trung
bình của chính cặp đó -- để hình phạt này có cùng bậc độ lớn với các thành
phần chi phí vận hành khác của cặp thay vì là một hằng số cố định dùng chung
cho mọi quy mô cầu; đây cũng chính là lý do thành phần này luôn dùng đơn giá
hằng số $c_{th}$ chứ không dùng $c_{th}^{(i)}$ theo giá thực: mục tiêu ``đạt
tối thiểu $85\%$ fill rate'' là một cam kết quản trị cố định, không nên tỷ lệ
thuận theo giá bán của từng mặt hàng.

\subsection{Chuẩn hóa phần thưởng theo từng cặp}
\label{subsec:chuanHoaReward}

Vì $300$ tác tử dùng chung một mạng nơ-ron (Mục~\ref{sec:kienTruc}), phần
thưởng thô ở công thức~(\ref{eq:localReward}) không thể đưa thẳng vào huấn
luyện: quy mô nhu cầu giữa các cặp trong bộ dữ liệu chênh lệch tới ba bậc độ
lớn (từ dưới $0{,}1$ đến hơn $150$ đơn vị/ngày), khiến phần thưởng thô của cặp
bán chạy nhất lớn hơn phần thưởng thô của cặp bán chậm nhất hàng nghìn lần.
Một mạng giá trị (critic) dùng chung trọng số cho toàn bộ $300$ cặp buộc phải
khớp đồng thời các mục tiêu trải trên nhiều bậc độ lớn như vậy, khiến sai số
huấn luyện bị chi phối gần như hoàn toàn bởi nhóm nhỏ các cặp có quy mô lớn
nhất, còn tín hiệu học của phần lớn các cặp còn lại gần như bị nhiễu lấn át.

Để khắc phục, phần thưởng cục bộ được chuẩn hóa theo một thang đo đặc trưng
của riêng từng cặp trước khi đưa vào bộ đệm huấn luyện:
\begin{equation}
    \tilde{r}_i^t = \frac{r_i^t}{c_{th}^{(i)} \cdot \bar{D}_i + c_{dh}}
    \label{eq:normReward}
\end{equation}
Mẫu số là chi phí ước lượng của ``một ngày xấu điển hình'' đối với cặp $i$:
hết hàng toàn bộ nhu cầu trong một ngày cộng thêm một lần phát sinh đơn đặt
hàng. Hằng số $c_{dh}$ ở mẫu số giữ cho phép chia không tiến về không đối với
những mặt hàng gần như không có giao dịch. Phép chuẩn hóa này chỉ đổi thang đo
của phần thưởng, không đổi thứ tự tốt/xấu giữa các hành động trong cùng một
cặp, nên không làm thay đổi chính sách tối ưu mà tác tử cần học; nó chỉ làm
cho gradient đóng góp bởi $300$ cặp có độ lớn tương đương nhau, cho phép mạng
dùng chung học đồng đều trên toàn bộ hệ thống thay vì chỉ tập trung vào những
cặp có quy mô chi phí lớn nhất. Mục~\ref{chap:chapter4} trình bày bằng chứng
thực nghiệm cho vai trò của cơ chế chuẩn hóa này.

\subsection{Vai trò của thành phần phạt fill rate}
\label{subsec:fillRate}

Thành phần $\phi_{dv} \cdot c_{th} \cdot \bar{D}_i \cdot \max(\tau -
\text{FR}_t^{(i, w)}, 0)$ trong công thức~(\ref{eq:localReward}) bổ sung một
tín hiệu phạt ràng buộc (constraint penalty) theo tinh thần Lagrangian: thay
vì chỉ phạt thiếu hàng tức thời trong ngày -- vốn có phương sai rất cao do đặc
tính gián đoạn của nhu cầu bán lẻ -- tỷ lệ đáp ứng tính trên cửa sổ trượt cung
cấp một đánh giá ổn định hơn về chất lượng phục vụ theo thời gian. Cần nói rõ
rằng thành phần phạt này \emph{chủ động thay đổi} nghiệm tối ưu của bài toán
tối ưu hóa chi phí thuần túy, nhằm điều hướng chính sách hướng tới việc thỏa
mãn một ràng buộc kinh doanh (mức phục vụ tối thiểu) mà bản thân hàm chi phí
vận hành không tự nhiên tạo ra động lực để đạt tới.

\subsection{Các tham số chi phí}
\label{subsec:thamSoChiPhi}

Bảng~\ref{tab:thamsoChiPhi} liệt kê các tham số chi phí sử dụng trong thực
nghiệm.

\begin{table}[htbp]
    \centering
    \caption{Tham số chi phí trong môi trường mô phỏng}
    \label{tab:thamsoChiPhi}
    \begin{tabular}{|l|c|c|}
        \hline
        \textbf{Thành phần} & \textbf{Ký hiệu} & \textbf{Giá trị} \\
        \hline
        Chi phí lưu kho      & $c_{lk}$    & $1{,}0$ /đơn vị/ngày \\ \hline
        Chi phí đặt hàng     & $c_{dh}$    & $5{,}0$ /lần đặt \\ \hline
        Phạt tràn kho        & $p_{tk}$    & $5{,}0$ /đơn vị tràn \\ \hline
        Phạt vi phạm mức phục vụ & $\phi_{dv}$ & $3{,}0$ \\ \hline
        Ngưỡng fill rate     & $\tau$      & $85\%$ \\ \hline
        Cửa sổ fill rate     & $w$         & $30$ ngày \\ \hline
        Đơn giá thiếu hàng dùng cho phạt mức phục vụ & $c_{th}$ & $10{,}0$ /đơn vị \\
        \hline
    \end{tabular}
\end{table}

Với $c_{th} = 10{,}0$ và $c_{lk} = 1{,}0$, tỷ số chi phí thiếu hàng ngầm định
mức phục vụ chu kỳ lý thuyết $\alpha \approx 90{,}9\%$ (Mục~\ref{subsec:thuocDoMucPhucVu}).
Do đó, ràng buộc fill rate $\beta \ge 85\%$ có khả năng rơi vào trạng thái
không có hiệu lực bắt buộc (non-binding) nếu chỉ xét riêng phần chi phí thiếu
hàng thuần túy; Chương~\ref{chap:chapter4} kiểm chứng lại giả thuyết này bằng
kết quả huấn luyện thực tế.

Ngoài đơn giá hằng số $c_{th}$ dùng cho thành phần phạt mức phục vụ, đơn giá
thiếu hàng thực tế $c_{th}^{(i)}$ dùng trong công thức~(\ref{eq:localReward})
và~(\ref{eq:normReward}) có thể được suy ra từ giá bán trung bình thực tế của
chính mặt hàng đó thay vì dùng một hằng số chung cho tất cả:
\begin{equation}
    c_{th}^{(i)} = \max\bigl(\bar{P}_i \cdot \rho,\; c_{th}^{\min}\bigr)
    \label{eq:realPrice}
\end{equation}
trong đó $\bar{P}_i$ là giá bán trung bình của cặp $i$ (ước lượng chỉ trên
miền huấn luyện, lấy từ dữ liệu giá bán theo tuần của M5), $\rho = 0{,}3$ là tỷ
lệ biên lợi nhuận giả định, và $c_{th}^{\min} = 0{,}5$ là mức sàn để tránh các
mặt hàng giá rất thấp có đơn giá thiếu hàng gần bằng không. \textbf{Đây là một
giả định phương pháp luận cần nêu rõ}: bộ dữ liệu M5 không công bố giá vốn của
từng mặt hàng, nên $\rho$ phản ánh giả định rằng một đơn vị thiếu hàng làm mất
đúng phần lợi nhuận biên của giao dịch đó, chứ không phải toàn bộ giá bán --
đây là một tham số cần trình bày tường minh như một giả định, không phải một
đại lượng đo được trực tiếp từ dữ liệu. Với cấu hình này, đơn giá thiếu hàng
suy ra trên $300$ cặp dao động trong khoảng $0{,}50$ đến $3{,}82$, trung bình
$1{,}20$ -- thấp hơn đáng kể so với hằng số $c_{th} = 10{,}0$ dùng cho thành
phần phạt mức phục vụ, đúng theo chủ đích thiết kế đã nêu ở Mục~\ref{subsec:localReward}.

\section{Kiến trúc mạng Actor--Critic chia sẻ tham số}
\label{sec:kienTruc}

\subsection{Hai mạng Actor và Critic độc lập, chia sẻ giữa các tác tử}
\label{subsec:sharedNet}

Toàn bộ $300$ tác tử dùng chung một bộ trọng số nơ-ron duy nhất
\parencite{gupta2017}: tại mỗi bước, quan sát của cả $300$ cặp được xếp thành
một batch và đi qua cùng một mạng trong một lượt tính toán duy nhất, thay vì
huấn luyện $300$ mạng riêng biệt. Nhờ đó, số tham số cần học là một hằng số,
không phụ thuộc vào quy mô hệ thống $N \times M$.

Điểm cần nhấn mạnh trong thiết kế là Actor (đầu ra phân phối xác suất trên
$K = 6$ mức đặt hàng) và Critic (đầu ra ước lượng giá trị trạng thái) là
\textbf{hai mạng hoàn toàn độc lập}, mỗi mạng gồm hai lớp ẩn kích thước $128$
với hàm kích hoạt $\tanh$ và khởi tạo trọng số trực giao, được tối ưu bằng hai
bộ tối ưu Adam riêng biệt và cắt chuẩn gradient (gradient clipping) riêng cho
từng mạng. Quyết định thiết kế này xuất phát từ một quan sát thực nghiệm quan
trọng trong quá trình phát triển: nếu Actor và Critic dùng chung một thân mạng
và một phép cắt gradient áp dụng cho tổng gradient của cả hai nhánh, độ lớn
gradient sinh ra từ hàm mất mát giá trị (value loss) có thể lớn hơn gradient
sinh ra từ hàm mất mát chính sách (policy loss) tới hàng trăm lần -- vì hai
hàm mất mát này có đơn vị và thang đo hoàn toàn khác nhau. Khi đó, phép cắt
gradient chung sẽ thu nhỏ toàn bộ gradient theo tỷ lệ do nhánh giá trị áp đảo
quyết định, khiến tốc độ học hiệu dụng của Actor bị giảm đi hàng chục đến hàng
trăm lần so với tốc độ học danh nghĩa, và chính sách gần như không nhận được
tín hiệu học trong giai đoạn đầu huấn luyện. Tách hai mạng và hai bước tối ưu
độc lập loại bỏ hoàn toàn sự giao thoa gradient này, trong khi vẫn giữ nguyên
luận điểm chia sẻ tham số \emph{giữa các tác tử} là không đổi.

\subsection{Ổn định hóa quá trình huấn luyện}
\label{subsec:onDinhHoa}

Bên cạnh việc tách hai mạng, ba cơ chế bổ sung được áp dụng để ổn định quá
trình huấn luyện trên quy mô $300$ tác tử:

\begin{itemize}
    \item \textbf{Chuẩn hóa mục tiêu giá trị (value normalization):} thay vì
    huấn luyện Critic khớp trực tiếp giá trị lợi tức thô -- vốn có thang đo
    thay đổi theo cấu hình chi phí -- lợi tức được chuẩn hóa bằng trung bình
    và phương sai trượt trước khi đưa vào hàm mất mát, và được khử chuẩn hóa
    khi trả giá trị ước lượng về cho quá trình thu thập dữ liệu. Cơ chế này
    tương tự kỹ thuật PopArt được sử dụng trong nhiều cài đặt MAPPO
    \parencite{yu2022mappo}, giữ hàm mất mát giá trị luôn ở một thang đo ổn
    định bất kể độ lớn tuyệt đối của chi phí vận hành.
    \item \textbf{Cắt giá trị (value clipping):} mỗi lần cập nhật, ước lượng
    giá trị mới bị giới hạn không lệch quá xa so với ước lượng cũ, tránh một
    lô dữ liệu bất thường kéo Critic đi quá xa trong một bước cập nhật.
    \item \textbf{Ràng buộc bước cập nhật chính sách:} huấn luyện được dừng
    sớm trong một vòng lặp cập nhật nếu độ phân kỳ Kullback--Leibler ước lượng
    giữa chính sách mới và chính sách cũ vượt quá một ngưỡng cho phép, theo
    đúng tinh thần vùng tin cậy của TRPO \parencite{schulman2015trpo} mà PPO
    kế thừa. Hệ số entropy và tốc độ học được giảm dần (annealing) theo tiến
    độ huấn luyện: hệ số entropy lớn ở giai đoạn đầu khuyến khích khám phá,
    sau đó giảm dần để chính sách hội tụ dứt khoát hơn ở giai đoạn cuối.
\end{itemize}

\subsection{Quy trình huấn luyện IPPO}
\label{subsec:quyTrinhTrain}

\begin{enumerate}
    \item Thu thập dữ liệu (rollout): chạy $n_{\text{steps}} = 1024$
    bước trong môi trường. Tại mỗi bước, cả $300$ tác tử cùng chọn hành động
    dựa trên quan sát cục bộ, tổng cộng thu được khoảng $1024 \times 300
    \approx 307.000$ mẫu $(o_i, a_i, r_i, o_i')$ mỗi vòng thu thập.

    \item Tính GAE: ước lượng lợi thế $\hat{A}_t^{(i)}$ cho từng tác
    tử $i$ \emph{độc lập theo trục thời gian} bằng công thức GAE
    (\ref{eq:gae}) với $\gamma = 0{,}99$ và $\lambda = 0{,}95$ -- đây chính là
    ý nghĩa chữ ``Independent'' trong IPPO: độc lập về cách tính lợi tức của
    từng tác tử, trong khi vẫn chia sẻ chung một bộ tham số mạng.

    \item Chuẩn hóa lợi thế theo từng cặp (per-pair advantage
    normalization): chuẩn hóa $\hat{A}_t^{(i)}$ theo trung bình và độ lệch
    chuẩn tính riêng cho mỗi cặp kho--mặt hàng $i$, cùng mục đích với phép
    chuẩn hóa phần thưởng ở Mục~\ref{subsec:chuanHoaReward}: giúp các cặp có
    quy mô chi phí khác nhau đóng góp cân bằng vào gradient của mạng dùng
    chung.

    \item Cập nhật PPO: thực hiện tối đa $E = 4$ epoch tối ưu hóa hàm mục
    tiêu clipped (\ref{eq:ppoclip}) trên toàn bộ dữ liệu rollout, chia thành
    các minibatch $16.384$ mẫu. Tốc độ học Actor $\eta_a = 3 \times 10^{-4}$,
    Critic $\eta_c = 1 \times 10^{-3}$, cả hai được giảm dần tuyến tính theo
    tiến độ huấn luyện.

    \item Lặp lại bước 1--4 trong $5000$ episode; cứ mỗi $25$ episode, chính
    sách được đánh giá theo phương thức xác định (không lấy mẫu ngẫu nhiên)
    trên miền kiểm định để chọn ra bộ trọng số tốt nhất, theo đúng phân chia
    miền dữ liệu ở Mục~\ref{subsec:phanChiaMien}.
\end{enumerate}

Bảng~\ref{tab:hyperparams} tóm tắt các siêu tham số chính.

\begin{table}[htbp]
    \centering
    \caption{Siêu tham số huấn luyện}
    \label{tab:hyperparams}
    \begin{tabular}{|l|l|c|}
        \hline
        \textbf{Nhóm} & \textbf{Tham số} & \textbf{Giá trị} \\
        \hline
        Môi trường & $N \times M$ (số tác tử) & $10 \times 30 = 300$ \\ \hline
                   & Episode length & $365$ ngày \\ \hline
                   & Lead time & $[1, 3]$ ngày \\ \hline
                   & Hệ số mức đặt hàng ($K$) & $\{0;\, 0{,}5;\, 1;\, 2;\, 3;\, 5\}$ \\
        \hline
        PPO & Rollout steps ($n_{\text{steps}}$) & $1024$ \\ \hline
            & PPO epochs ($E$) & $4$ \\ \hline
            & Kích thước minibatch & $16.384$ \\ \hline
            & Learning rate Actor & $3 \times 10^{-4}$ \\ \hline
            & Learning rate Critic & $1 \times 10^{-3}$ \\ \hline
            & $\gamma$ & $0{,}99$ \\ \hline
            & $\lambda_{\text{GAE}}$ & $0{,}95$ \\ \hline
            & Clip $\epsilon$ & $0{,}2$ \\ \hline
            & Hệ số entropy (đầu $\to$ cuối) & $0{,}02 \to 0{,}002$ \\ \hline
            & Ngưỡng KL cho dừng sớm & $0{,}02$ \\ \hline
            & Total episodes & $5000$ \\
        \hline
        Mạng & Kiến trúc & $2$ mạng độc lập, mỗi mạng $2 \times 128$ \\ \hline
              & Activation & $\tanh$ \\
        \hline
    \end{tabular}
\end{table}

\section{Thiết kế môi trường mô phỏng}
\label{sec:moiTruong}

Môi trường được xây dựng theo chuẩn giao diện Gymnasium
\parencite{towers2024gymnasium}, kế thừa từ chuẩn giao diện OpenAI Gym
\parencite{brockman2016gym}: mỗi tập mô phỏng có thể được khởi tạo lại một
cách độc lập và tái lập được nếu cùng một hạt giống ngẫu nhiên, và tiến triển
theo từng bước thời gian rời rạc dựa trên hành động đồng thời của toàn bộ tác
tử. Môi trường quản lý ma trận tồn kho của $300$ cặp kho--mặt hàng, hàng đợi
giao hàng đang trên đường theo từng cặp, lịch sử nhu cầu và tỷ lệ đáp ứng
trượt, cùng hàm tính chi phí và phần thưởng theo công thức~(\ref{eq:localReward}).
Mỗi episode kéo dài $365$ bước, tương đương một năm kinh doanh.

Khóa luận chủ động tự xây dựng môi trường mô phỏng này thay vì sử dụng trực
tiếp các bộ thư viện benchmark sẵn có như OR-Gym \parencite{hubbs2020orgym}
hay MABIM \parencite{yang2023mabim}. Lý do cốt lõi là các benchmark hiện hành
chưa hỗ trợ sẵn đồng thời cấu hình mạng lưới $10$ kho song song độc lập nguồn
cung, ràng buộc sức chứa kho hiệu chỉnh riêng theo từng kho và dùng chung giữa
$30$ SKU, cùng cơ chế phạt mức phục vụ theo cửa sổ trượt như mô tả ở
Mục~\ref{sec:hamThuong}. Việc tự phát triển môi trường giúp đảm bảo khả năng
kiểm soát toàn diện các tham số mô phỏng và tính tương thích cao nhất với mục
tiêu nghiên cứu của đề tài.
